\chapter{第七讲\quad 客观要件三：因果关系}
\label{lecture:07}
\noindent{\small\hyperref[ref:lecture:07]{查阅本讲扩展规则}}\par

对故意结果犯，归责关系关系到法定结果是否算作既遂结果。对要求实害的过失犯罪，关系到能否成立该罪。对结果加重犯，关系到能否升档。把三个问题分开，依次记录「自然过程是否连得上」「结果是否落入该规范」「证据能否证明」。罪过另见第八讲。

\section{第一节\quad 关于“因”的要求}

\paragraph{定义与教材模型}
令 $a$ 为特定行为，$g$ 为受保护法益，$r$ 为具体结果，$j$ 为所审罪名。$\mathrm{New}(a,g)$ 表示行为创设或增加了与该结果有关的风险，$\mathrm{Perm}_j(a,g)$ 表示该风险在所审规范下仍获允许，$\mathrm{Risk}_j(a,g)$ 表示该风险超出允许范围。教材先筛出
\[
\mathrm{Risk}_j(a,g)=\mathrm{New}(a,g)\land\neg\mathrm{Perm}_j(a,g).
\]
过失型案件须指明何种注意义务被违反，查清义务所针对的危险。合法医疗、按规范进行交通等活动本身可能存在风险，其风险是否获允许要依据具体规则和当时事实。可允许风险和正当防卫等违法阻却理由所属的判断层次可以随理论体系不同而变动。教材的考查路径侧重排除未经行为创设或未超出允许范围的风险。

\paragraph{推导}
驾驶人按当时视野、速度等条件遵守交通义务，未获任何行人突然越过隔离带的预警，题设已证明正常操作亦无法避免碰撞。在此限定事实下，车辆运动是撞击的物理条件，记作 $C_{\rm phys}=1$。已证明待审风险属于允许的通常交通风险，$\mathrm{Perm}_j=1$，故
\[
C_{\rm phys}=1,\quad
\mathrm{Risk}_j=\mathrm{New}\land 0=0.
\]
对于依违反注意义务来审查的过失型罪名，不能仅凭碰撞和死亡得出刑法归责。若甲仅希望乙在平常跑步中遇险，却未实施制造风险的行为，则更早在 $\mathrm{New}$ 一项即停止。

\paragraph{边界}
信赖原则限于合理信赖。若驾驶人已看见道路中的儿童，仍高速向前行驶，不能继续把「他人本应守规则」当作 $\mathrm{Perm}_j$ 的充分证据。用事故前可见情形评估风险和义务，不能用事故后的严重程度回填。

\section{第二节\quad 关于“果”的要求}

\paragraph{定义与教材模型}
令 $r=(v,h,t,\mu)$ 表示被害对象 $v$、损害种类 $h$、实际发生时段 $t$ 与可证的机制 $\mu$。记录 $\mathrm{Real}(r)$ 表示实害已经发生，$\mathrm{Scope}_j(a,r)$ 表示它落入罪状及有关注意规范保护范围，$\mathrm{Sphere}(a,r)$ 表示它属于行为人负责控制的风险领域。可归责的结果须先通过
\[
\mathrm{Result}_j(a,r)
   =\mathrm{Real}(r)\land\mathrm{Scope}_j(a,r)
      \land\mathrm{Sphere}(a,r).
\]
本处论实害归责。危险犯的法定危险作为独立构成条件，依第六讲审查。对注意义务的保护目的，应给出所采规范及结果类型，不能只背典型题的结论。

\paragraph{推导}
甲违规碰撞乙，导致乙昏迷。与甲无联络的丙随后独立盗走乙的钱包。令 $r_1$ 为撞击伤害，$r_2$ 为钱包被盗造成的财产损失。题设给出 $\mathrm{Real}(r_1)\land\mathrm{Real}(r_2)$。即使反事实模型认为没有碰撞就没有此次盗取，碰撞安全义务所防止的直接身体伤害或碰撞损失与独立盗取的财产损失不同。因此
\[
\mathrm{Scope}_j(a,r_2)=0
  \ \Longrightarrow\ \mathrm{Result}_j(a,r_2)=0 .
\]
碰撞是否另构成交通肇事罪、盗取是否构成盗窃等财产罪，须分别满足各自条件。仅有甲先前违法，不能把丙的损失结果转给甲。

\paragraph{真实结果、假设后备原因与并行风险}
如果甲在行刑官依法执行前擅自造成受刑人死亡，则待解释的是甲在实际时点造成的死亡事件。日后可能由行刑官合法造成另一个死亡事件，只是后备原因。若后行为阻断了前行为的危险流，例如毒水未被饮用而被放空，被害人实际因无水而渴死，应分别找实际经过的机制和从未实现的机制。

取两个独立投放行为，实际均发挥作用。为区分事实机制，构造简化阈值模型
\[
R(x_1,x_2)=\mathbf{1}_{\{d_1x_1+d_2x_2\ge 10\}},
\qquad x_i\in\{0,1\}.
\]
这里 $x_i$ 表示第 $i$ 次投放是否存在，$d_i$ 是剂量，十是题设阈值。医学和法律均不从此人为数值模型推出通用阈值。

\begin{enumerate}
\item 当 $(d_1,d_2)=(10,10)$，有 $R(1,1)=R(1,0)=R(0,1)=1$，同时 $R(0,0)=0$。逐一删除任一投放都不能使死亡消失，但每份剂量单独足够，两个危险在实际均起作用。教材称竞合的因果关系，二因一果。
\item 当 $(d_1,d_2)=(5,5)$，有 $R(1,1)=1$，$R(1,0)=R(0,1)=0$。每次投放都是该组合的必要条件，单独都不充分。若采教材对独立故意行为的单独危险实现说，且两人没有共同犯罪关系，则各自只能就自己的投放负责，题设符合未遂要件时分别论未遂。若采实际机制的累积贡献归因，二者均可被视为事实原因，再另论法律归责。这里变的是理论规则，四个计算值相同。法律归责不能直接从剂量计算推出。
\end{enumerate}

\paragraph{边界}
警察接管车辆检查后，因其独立控制中的新风险造成追尾，教材要求另查 $\mathrm{Sphere}$。职责可能并存，接管本身不必然消除原驾驶人的危险。若原驾驶人仍可控制或持续制造同一风险，应保留双重归责的可能，再查实际机制。

\section{第三节\quad 关于“因”与“果”的关系：介入因素的案件}

\paragraph{反事实与机制}
在确定的模型 $M$ 和背景 $u$ 中，用 $R_{\operatorname{do}(a=0)}(u)$ 表示取消 $a$ 后，待审结果是否发生。干预保持背景条件，让行为的后续变量按模型重新计算。若实际 $R(u)=1$ 且 $R_{\operatorname{do}(a=0)}(u)=0$，则 $a$ 是该背景下结果的反事实必要条件。对充分性也须给出背景。例如固定在 $u$ 中并指定共存条件，才可讨论 $R_{\operatorname{do}(a=1)}(u)=1$。两者都不是任意场景的蕴含。并行充分原因可以使反事实必要条件检验遗漏事实作用，纯粹引人到另一独立风险所在也可以通过反事实检验。对后者仍须进行法律归责。

\paragraph{危险现实化的教材规则}
令 $a$ 为先前行为，$i$ 为后续因素。先查 $i$ 是否在具体情境中通常由 $a$ 引发。记该引发关系为 $T(a,i)$。若 $T(a,i)$ 已证且 $r$ 确为该引发风险所实现，在规范保护范围内可以将 $r$ 归属给 $a$。对 $i$ 的归责仍单独判断。若 $i$ 异常而相对独立，再查 $r$ 的实际机制是前危险占主导、后危险占主导，还是两种危险均实质参与。用 $\mathrm{Dom}_a,\mathrm{Dom}_i,\mathrm{Both}$ 记录教材所示三种选择。这里的「作用」是证据支持的机制与规范评价，不是主观打分或时间先后。
\[
\begin{array}{c|c}
\text{已证明的具体机制} & \text{该模型允许的归属}\\ \hline
T(a,i)\text{ 且引发危险实现} & a\text{，并另审 }i\\
\text{独立介入且 }\mathrm{Dom}_a & a\\
\text{独立介入且 }\mathrm{Dom}_i & i\\
\text{两种危险均实质实现 }\mathrm{Both} & a,\ i
\end{array}
\]
每格还以创设不允许风险、保护范围等前提为限。证据不足以比较作用时，转第四节。

\paragraph{推导}
甲持刀袭击乙，乙负伤送医。途中，第三人丙严重违规撞车，乙仅因撞击死亡。令 $a$ 为袭击，$i$ 为独立交通事故。撤去 $a$ 可能没有此次行车，故反事实检验可能通过。题设进一步给出
\[
\Gamma\vdash
  \neg T(a,i)\land\mathrm{Dom}_i
   \land\mathrm{Risk}_j(i,g)\land\mathrm{Scope}_j(i,r).
\]
根据第二步，仅将这次死亡归属于撞击危险。甲按已证的杀人故意和着手等条件审查未遂，若其故意仅伤害则按伤害罪要件审查。丙是否犯罪还须查罪过和交通肇事罪门槛。若改成甲把乙推至高速车道，紧随车辆来不及躲避而碾轧，车流危险为先前行为所引发，$T(a,i)$ 可能为真，此时死亡可归给前行为。

\paragraph{边界与四种介入}
对被害人、第三人、救助中断和特殊体质应逐类找机制。逃避眼前袭击而跌落，和有自由选择时事后自杀的情境不同。对持续虐待或暴力干涉婚姻自由导致的自杀，教材要求审查累积胁迫和结果加重条文，不能套一般自杀结论。轻微救治偏差并不必然切断重伤至死亡的机制。异常错误的治疗若独自决定死亡，结论才可能反转。现实且有实效的救助已经展开，后来被他人故意阻断时，应查阻断者是否决定了死亡。尚无现实救助，不能用「本来有人会救」虚构同一结构。

特殊体质是行为前的背景，发病或病情恶化才是后续事件。若伤害使患者原有凝血障碍转化为持续出血死亡，可认定伤害与死亡有客观因果联系。行为人能否预见死亡另审。若无死亡罪过，不能只靠因果关系成立死亡加重档。

\section{第四节\quad 无法查明的案件}

\paragraph{证据模型与法条}
《刑事诉讼法》第五十五条要求定罪量刑的事实有证据证明，综合全案证据对所认定事实排除合理怀疑。第二百条规定证据不足而不能认定有罪时的处理。为训练这个约束，令 $\mathcal H_\Gamma$ 是有证据根据、仍未合理排除的事实情形，令 $E_x$ 为对被告 $x$ 不利的具体要素，且 $\mathcal H_\Gamma\ne\varnothing$。课堂上的充分核验式是
\[
\forall H\in\mathcal H_\Gamma,\quad H\vdash E_x.
\]
每个 $H$ 中使用同一选定法律规则分别推导。事实缺口不能写成 $\neg E_x$ 已证。若仍有合理情形中该要素未成立，则不能把该要素按成立处理。结果轻重按各罪的全部法定条件审查，死亡越早并非普遍的量刑排序。

\paragraph{推导一，两人无共同犯罪}
甲乙各自独立，均有杀丙故意且已经射击。头部中弹致死，腿部中弹无独立致死作用。身份鉴定不能查明哪颗子弹出自谁。令 $E_x$ 表示 $x$ 的个人行为导致死亡，$\Gamma$ 已分别证明二人均有杀人故意和着手。两种合理情形为
\[
\begin{array}{c|cc}
 & E_{\text{甲}} & E_{\text{乙}}\\ \hline
H_1\colon\text{甲弹击中头部} &1&0\\
H_2\colon\text{乙弹击中头部} &0&1
\end{array}
\]
对甲，$H_2$ 使致死要素不能确定。对乙，$H_1$ 起同样作用。若没有其他结果归责规则，在分别满足未遂条件的题设下按未遂处理。判定依据是对每人不能证明致死归责，不能在事实叙述中同时认定「两人都未打中头部」。

\paragraph{推导二，存在共同正犯}
改为已证明二人共谋且共同实行，死亡属于共同故意和共同实行范围。令 $C_{\mathrm{joint}}$ 为这一共同归责前提，$G_x$ 表示共同结果能归给共同正犯 $x$。教材规则为
\[
C_{\mathrm{joint}}\land(E_{\text{甲}}\lor E_{\text{乙}})
 \ \Longrightarrow\ G_{\text{甲}}\land G_{\text{乙}}.
\]
两种合理情形均有致命一枪来自共同实行，所以经分类讨论均推出 $G_{\text{甲}}\land G_{\text{乙}}$。无需证明每人的单独射击均物理致命，也不能以此省略共同犯罪事实的证明。

\paragraph{边界}
在连续碾轧题中，分别设「第一次后已死亡」「第二次后才死亡」两种情形。若第一人使被害人倒在车流中，第二次碾轧是其危险的通常实现，则前一人可在两情形中都被归责，第二人只在后一情形可归责。若题设改成第一次后被害人已脱离危险，后来被独立运到另一车道，应重新推导两个分支。类似地，取走财物时所有人是否死亡只是一个事实变量，是否存在他人占有、侵占的数额与拒不退还等条件仍须证明。

\section*{闭卷记忆与自测}

\paragraph{复述路线}
先背「因、果、链、证、心」。因写新风险和允许界线。果写现实、规范目的和管辖。链画前行为、介入、结果，再写「通常引发，独立比作用」。证用两行表逐人检查合理情形。心单独写预见、希望、放任或过失，避免把归责直接等同有责。

\paragraph{题一}
分别算出两份剂量 $(10,10)$ 与 $(5,5)$ 时，逐一删除投放的四个值。为何同样见到死亡，不能直接套同一个条件判断？

\paragraph{核对}
前者删除任一份仍为一。后者删除任一份即为零。前者各份单独充分，后者各份仅在共同机制中是必要的。事实机制与教材采用的归责规则应分开写。

\paragraph{题二}
受伤后昏迷使患者吸入异物窒息，与获救出院时遭遇独立纵火而死，哪一项应先比较？

\paragraph{核对}
先查 $T(a,i)$。昏迷引发的异物吸入可能是伤害危险的实现。独立纵火须进一步比对实际机制，若已独立决定死亡，先前伤害不能仅因提供时空机会而归属该死亡。

\paragraph{题三}
互不联络的两人各打一枪，只能证明其中一枪致命。结论可以写「两人都没有致死」吗？

\paragraph{核对}
不可以。共同事实中有一人致死，个体致死要素各自不能排除合理怀疑。写出对每名被告的证明结论，不制造一个与已证事实矛盾的故事。
